% Traduction française COMIMa de l'énoncé original (anglais).
\emph{Suite de Morse–Thue.} On part de $0$ ; à chaque segment initial, on ajoute son complémentaire : $0,\ 01,\ 0110,\ 01101001, \ldots$.

(a) On note $x(0), x(1), x(2), \ldots$ les chiffres de la suite. Montrer que $x(2n) = x(n)$ et $x(2n+1) = 1 - x(2n)$.

(b) Montrer que $x(n) = 1 - x(n - 2^k)$, où $2^k$ est la plus grande puissance de $2$ inférieure ou égale à $n$. Trouver le $1993$\textsuperscript{e} chiffre de la suite.

(c) Montrer que la suite n'est pas périodique.

(d) On écrit les entiers positifs ou nuls en base $2$ : $0, 1, 10, 11, \ldots$, puis on remplace chacun par la somme de ses chiffres modulo $2$. On obtient la suite de Morse–Thue. Le démontrer.
